% ------------------------------------------------------------------------------
% Slide 4: El Dilema de la Vocal I: Sonrisa - Correlación vs Lengua - Causa
% ------------------------------------------------------------------------------
\begin{frame}{El Dilema de la Vocal I: Sonrisa - Correlación vs Lengua - Causa}
    \begin{columns}[c]
        \begin{column}{0.51\textwidth}
            \scriptsize
            \begin{alertblock}{La Sonrisa en Mejilla no es la Causa Acústica}
                \begin{itemize}\setlength{\itemsep}{1pt}\setlength{\parskip}{0pt}
                    \item Al decir la /i/ se suele sonreír (\textit{cigomático}), pero es un gesto que solo acompaña.
                    \item La causa física real que sube $F_2$ es \textbf{empujar la lengua hacia adelante y arriba}.
                    \item Se puede pronunciar una /i/ limpia sin mover los labios, usando solo la lengua.
                \end{itemize}
            \end{alertblock}
            \vspace{0.1cm}
            \begin{exampleblock}{Qué Pasó en las Mediciones a lo Largo del Tiempo}
                Con los 3 sensores clásicos (\textit{Cigomático, Digástrico, Orbicular}):
                \begin{itemize}\setlength{\itemsep}{1pt}\setlength{\parskip}{0pt}
                    \item \textbf{Al principio:} Alta separación porque la sonrisa estaba acoplada a la /i/.
                    \item \textbf{Meses después:} El hablante sonrió también en la /a/, o hizo la /i/ solo con la lengua sin mover la cara, perdiendo separación.
                \end{itemize}
            \end{exampleblock}
        \end{column}
        \begin{column}{0.47\textwidth}
            \centering
            \placeholderfigura{0.70\textheight}{Comparativa de PCA del Sujeto 1: Antes - alta separabilidad vs. Después - colapso entre A, E, I por desacople del gesto de sonrisa}
        \end{column}
    \end{columns}
\end{frame}
